%% Generated by Sphinx.
\def\sphinxdocclass{report}
\documentclass[letterpaper,10pt,english]{sphinxmanual}
\ifdefined\pdfpxdimen
   \let\sphinxpxdimen\pdfpxdimen\else\newdimen\sphinxpxdimen
\fi \sphinxpxdimen=.75bp\relax
\ifdefined\pdfimageresolution
    \pdfimageresolution= \numexpr \dimexpr1in\relax/\sphinxpxdimen\relax
\fi
%% let collapsible pdf bookmarks panel have high depth per default
\PassOptionsToPackage{bookmarksdepth=5}{hyperref}

\PassOptionsToPackage{booktabs}{sphinx}
\PassOptionsToPackage{colorrows}{sphinx}

\PassOptionsToPackage{warn}{textcomp}
\usepackage[utf8]{inputenc}
\ifdefined\DeclareUnicodeCharacter
% support both utf8 and utf8x syntaxes
  \ifdefined\DeclareUnicodeCharacterAsOptional
    \def\sphinxDUC#1{\DeclareUnicodeCharacter{"#1}}
  \else
    \let\sphinxDUC\DeclareUnicodeCharacter
  \fi
  \sphinxDUC{00A0}{\nobreakspace}
  \sphinxDUC{2500}{\sphinxunichar{2500}}
  \sphinxDUC{2502}{\sphinxunichar{2502}}
  \sphinxDUC{2514}{\sphinxunichar{2514}}
  \sphinxDUC{251C}{\sphinxunichar{251C}}
  \sphinxDUC{2572}{\textbackslash}
\fi
\usepackage{cmap}
\usepackage[T1]{fontenc}
\usepackage{amsmath,amssymb,amstext}
\usepackage{babel}



\usepackage{tgtermes}
\usepackage{tgheros}
\renewcommand{\ttdefault}{txtt}



\usepackage[Bjarne]{fncychap}
\usepackage{sphinx}

\fvset{fontsize=auto}
\usepackage{geometry}


% Include hyperref last.
\usepackage{hyperref}
% Fix anchor placement for figures with captions.
\usepackage{hypcap}% it must be loaded after hyperref.
% Set up styles of URL: it should be placed after hyperref.
\urlstyle{same}

\addto\captionsenglish{\renewcommand{\contentsname}{Getting Started:}}

\usepackage{sphinxmessages}
\setcounter{tocdepth}{3}
\setcounter{secnumdepth}{3}


\title{Stochpy Documentation}
\date{Dec 30, 2023}
\release{1.0.0}
\author{Alex Villarroel}
\newcommand{\sphinxlogo}{\vbox{}}
\renewcommand{\releasename}{Release}
\makeindex
\begin{document}

\ifdefined\shorthandoff
  \ifnum\catcode`\=\string=\active\shorthandoff{=}\fi
  \ifnum\catcode`\"=\active\shorthandoff{"}\fi
\fi

\pagestyle{empty}
\sphinxmaketitle
\pagestyle{plain}
\sphinxtableofcontents
\pagestyle{normal}
\phantomsection\label{\detokenize{index::doc}}


\sphinxstepscope


\chapter{Slip\_generation}
\label{\detokenize{modules:slip-generation}}\label{\detokenize{modules::doc}}
\sphinxstepscope


\section{Main package}
\label{\detokenize{main:main-package}}\label{\detokenize{main::doc}}

\subsection{Module contents}
\label{\detokenize{main:module-main}}\label{\detokenize{main:module-contents}}\index{module@\spxentry{module}!main@\spxentry{main}}\index{main@\spxentry{main}!module@\spxentry{module}}\begin{description}
\sphinxlineitem{Module with functions necessary for the creation of faults and the calculation of slip distributions}
\sphinxAtStartPar
Author/Github: Alex Villarroel Carrasco / @alexvillarroel
University of Concepcion,Chile.

\sphinxAtStartPar
SOME FUNCTIONS ARE COPYRIGHTED BY RODRIGO CIFUENTES LOBOS
in his github repository: @RCifuentesLobos : PaleoTsunami

\sphinxAtStartPar
/

\end{description}
\index{Mw\_kauselramirez() (in module main)@\spxentry{Mw\_kauselramirez()}\spxextra{in module main}}

\begin{fulllineitems}
\phantomsection\label{\detokenize{main:main.Mw_kauselramirez}}
\pysigstartsignatures
\pysiglinewithargsret{\sphinxcode{\sphinxupquote{main.}}\sphinxbfcode{\sphinxupquote{Mw\_kauselramirez}}}{\sphinxparam{\DUrole{n}{ms}}}{}
\pysigstopsignatures
\end{fulllineitems}

\index{N\_max\_matriz\_covarianza() (in module main)@\spxentry{N\_max\_matriz\_covarianza()}\spxextra{in module main}}

\begin{fulllineitems}
\phantomsection\label{\detokenize{main:main.N_max_matriz_covarianza}}
\pysigstartsignatures
\pysiglinewithargsret{\sphinxcode{\sphinxupquote{main.}}\sphinxbfcode{\sphinxupquote{N\_max\_matriz\_covarianza}}}{\sphinxparam{\DUrole{n}{C}}}{}
\pysigstopsignatures
\sphinxAtStartPar
Busca la esquina de la curva l de los valores propios
Entrada:
C : matriz de covarianza de la falla

\end{fulllineitems}

\index{calcula\_laplaciano() (in module main)@\spxentry{calcula\_laplaciano()}\spxextra{in module main}}

\begin{fulllineitems}
\phantomsection\label{\detokenize{main:main.calcula_laplaciano}}
\pysigstartsignatures
\pysiglinewithargsret{\sphinxcode{\sphinxupquote{main.}}\sphinxbfcode{\sphinxupquote{calcula\_laplaciano}}}{\sphinxparam{\DUrole{n}{tamano\_sf}}\sphinxparamcomma \sphinxparam{\DUrole{n}{F}}}{}
\pysigstopsignatures
\sphinxAtStartPar
Entradas 
tamano\_sf: tamano subfalla
F: Campo escalar a calcularle el laplaciano

\end{fulllineitems}

\index{corr2d\_fourier() (in module main)@\spxentry{corr2d\_fourier()}\spxextra{in module main}}

\begin{fulllineitems}
\phantomsection\label{\detokenize{main:main.corr2d_fourier}}
\pysigstartsignatures
\pysiglinewithargsret{\sphinxcode{\sphinxupquote{main.}}\sphinxbfcode{\sphinxupquote{corr2d\_fourier}}}{\sphinxparam{\DUrole{n}{slip}}\sphinxparamcomma \sphinxparam{\DUrole{n}{test}}\sphinxparamcomma \sphinxparam{\DUrole{n}{lonslip}\DUrole{o}{=}\DUrole{default_value}{{[}{]}}}\sphinxparamcomma \sphinxparam{\DUrole{n}{latslip}\DUrole{o}{=}\DUrole{default_value}{{[}{]}}}\sphinxparamcomma \sphinxparam{\DUrole{n}{lontest}\DUrole{o}{=}\DUrole{default_value}{{[}{]}}}\sphinxparamcomma \sphinxparam{\DUrole{n}{lattest}\DUrole{o}{=}\DUrole{default_value}{{[}{]}}}\sphinxparamcomma \sphinxparam{\DUrole{n}{normalizar}\DUrole{o}{=}\DUrole{default_value}{True}}}{}
\pysigstopsignatures
\sphinxAtStartPar
Calcula la correlacion cruzada entre 2 arrays 2d utilizando el metodo de la 
transformada rapida de fourier para obtener indices para filtrar
Entradas:
slip: array que contiene el slip generado
test: arrray de prueba con el cual se busca realzar el filtro (e.g.: anom. de gravedad)
lonslip, latslip: longitud y latitud de las subfallas que componen el slip
lontest, lattest: longitud y latitud de las subfallas que componen el test

\end{fulllineitems}

\index{crea\_falla() (in module main)@\spxentry{crea\_falla()}\spxextra{in module main}}

\begin{fulllineitems}
\phantomsection\label{\detokenize{main:main.crea_falla}}
\pysigstartsignatures
\pysiglinewithargsret{\sphinxcode{\sphinxupquote{main.}}\sphinxbfcode{\sphinxupquote{crea\_falla}}}{\sphinxparam{\DUrole{n}{lats}}\sphinxparamcomma \sphinxparam{\DUrole{n}{lons}}\sphinxparamcomma \sphinxparam{\DUrole{n}{prof}}\sphinxparamcomma \sphinxparam{\DUrole{n}{dip}}\sphinxparamcomma \sphinxparam{\DUrole{n}{strike}}\sphinxparamcomma \sphinxparam{\DUrole{n}{latini}}\sphinxparamcomma \sphinxparam{\DUrole{n}{latfin}}\sphinxparamcomma \sphinxparam{\DUrole{n}{area\_sf}}\sphinxparamcomma \sphinxparam{\DUrole{n}{profundidad}}\sphinxparamcomma \sphinxparam{\DUrole{n}{razon\_aspecto}}}{}
\pysigstopsignatures
\sphinxAtStartPar
Crea una falla a partir de datos del modelo Slab2.0 de Hayes, 2018
Entradas:
lats: latitudes de Slab2.0
lons: longitudes de Slab2.0
prof: profundidades de Slab2.0
dip: dip de Slab2.0
strike: strike de Slab2.0
latini: latitud inicial (norte)
latfin: latitud final (sur)
area\_sf: Area deseada para cada subfalla
profundidad: profundidad media de la falla, sobre la cual se calculara el punto medio de la distribucion
razon\_aspecto : Razon de aspecto de la falla (largo/ancho) se calcula el ancho de la falla a partir del largo, para respetar la relacion L/A = const. de Kanamori Anderson, 1975

\end{fulllineitems}

\index{deg2km() (in module main)@\spxentry{deg2km()}\spxextra{in module main}}

\begin{fulllineitems}
\phantomsection\label{\detokenize{main:main.deg2km}}
\pysigstartsignatures
\pysiglinewithargsret{\sphinxcode{\sphinxupquote{main.}}\sphinxbfcode{\sphinxupquote{deg2km}}}{\sphinxparam{\DUrole{n}{deg}}}{}
\pysigstopsignatures
\end{fulllineitems}

\index{dist\_haversine() (in module main)@\spxentry{dist\_haversine()}\spxextra{in module main}}

\begin{fulllineitems}
\phantomsection\label{\detokenize{main:main.dist_haversine}}
\pysigstartsignatures
\pysiglinewithargsret{\sphinxcode{\sphinxupquote{main.}}\sphinxbfcode{\sphinxupquote{dist\_haversine}}}{\sphinxparam{\DUrole{n}{lon1}}\sphinxparamcomma \sphinxparam{\DUrole{n}{lon2}}\sphinxparamcomma \sphinxparam{\DUrole{n}{lat1}}\sphinxparamcomma \sphinxparam{\DUrole{n}{lat2}}}{}
\pysigstopsignatures
\end{fulllineitems}

\index{dist\_sf() (in module main)@\spxentry{dist\_sf()}\spxextra{in module main}}

\begin{fulllineitems}
\phantomsection\label{\detokenize{main:main.dist_sf}}
\pysigstartsignatures
\pysiglinewithargsret{\sphinxcode{\sphinxupquote{main.}}\sphinxbfcode{\sphinxupquote{dist\_sf}}}{\sphinxparam{\DUrole{n}{lon1}}\sphinxparamcomma \sphinxparam{\DUrole{n}{lon2}}\sphinxparamcomma \sphinxparam{\DUrole{n}{lat1}}\sphinxparamcomma \sphinxparam{\DUrole{n}{lat2}}}{}
\pysigstopsignatures
\sphinxAtStartPar
Calcula la distancia en metros entre dos subfallas (sf) i y j utilizando el metodo de karney, 2013 y el elipsoide wgs84
Entradas:
lon1: longitud de subfalla i\sphinxhyphen{}esima
lon2: longitud de subfalla j\sphinxhyphen{}esima
lat1: latitud de subfalla i\sphinxhyphen{}esima
lat2: latitud de subfalla j\sphinxhyphen{}esima

\end{fulllineitems}

\index{dist\_sf\_alt() (in module main)@\spxentry{dist\_sf\_alt()}\spxextra{in module main}}

\begin{fulllineitems}
\phantomsection\label{\detokenize{main:main.dist_sf_alt}}
\pysigstartsignatures
\pysiglinewithargsret{\sphinxcode{\sphinxupquote{main.}}\sphinxbfcode{\sphinxupquote{dist\_sf\_alt}}}{\sphinxparam{\DUrole{n}{lon1}}\sphinxparamcomma \sphinxparam{\DUrole{n}{lon2}}\sphinxparamcomma \sphinxparam{\DUrole{n}{lat1}}\sphinxparamcomma \sphinxparam{\DUrole{n}{lat2}}}{}
\pysigstopsignatures
\sphinxAtStartPar
Calcula la distancia en metros entre dos subfallas (sf) i y j resolviendo el problema inverso 
de la geodesia. Utiliza el elipsoide WGS84. 
Entradas:
lon1: longitud de subfalla i\sphinxhyphen{}esima
lon2: longitud de subfalla j\sphinxhyphen{}esima
lat1: latitud de subfalla i\sphinxhyphen{}esima
lat2: latitud de subfalla j\sphinxhyphen{}esima

\end{fulllineitems}

\index{dist\_sf\_vectorized() (in module main)@\spxentry{dist\_sf\_vectorized()}\spxextra{in module main}}

\begin{fulllineitems}
\phantomsection\label{\detokenize{main:main.dist_sf_vectorized}}
\pysigstartsignatures
\pysiglinewithargsret{\sphinxcode{\sphinxupquote{main.}}\sphinxbfcode{\sphinxupquote{dist\_sf\_vectorized}}}{\sphinxparam{\DUrole{n}{lon1}}\sphinxparamcomma \sphinxparam{\DUrole{n}{lon2}}\sphinxparamcomma \sphinxparam{\DUrole{n}{lat1}}\sphinxparamcomma \sphinxparam{\DUrole{n}{lat2}}}{}
\pysigstopsignatures
\end{fulllineitems}

\index{distribucion\_slip() (in module main)@\spxentry{distribucion\_slip()}\spxextra{in module main}}

\begin{fulllineitems}
\phantomsection\label{\detokenize{main:main.distribucion_slip}}
\pysigstartsignatures
\pysiglinewithargsret{\sphinxcode{\sphinxupquote{main.}}\sphinxbfcode{\sphinxupquote{distribucion\_slip}}}{\sphinxparam{\DUrole{n}{C}}\sphinxparamcomma \sphinxparam{\DUrole{n}{mu}}\sphinxparamcomma \sphinxparam{\DUrole{n}{N}}}{}
\pysigstopsignatures
\sphinxAtStartPar
Calcula la distribucion de slip con la expansion de karhunen\sphinxhyphen{}loeve log\sphinxhyphen{}normal
exp(mu)exp(sum(zk*sqrt(lambdak)*vk))
Entradas: 
C: matriz de covarianza, se le calculan los eigenvalues y vectors
mu: matriz con medias
N: numero de modos a contar para la sumatoria

\end{fulllineitems}

\index{escalar\_magnitud\_momento() (in module main)@\spxentry{escalar\_magnitud\_momento()}\spxextra{in module main}}

\begin{fulllineitems}
\phantomsection\label{\detokenize{main:main.escalar_magnitud_momento}}
\pysigstartsignatures
\pysiglinewithargsret{\sphinxcode{\sphinxupquote{main.}}\sphinxbfcode{\sphinxupquote{escalar\_magnitud\_momento}}}{\sphinxparam{\DUrole{n}{Mw}}\sphinxparamcomma \sphinxparam{\DUrole{n}{slip}}\sphinxparamcomma \sphinxparam{\DUrole{n}{prof}}\sphinxparamcomma \sphinxparam{\DUrole{n}{lons}}\sphinxparamcomma \sphinxparam{\DUrole{n}{lats}}\sphinxparamcomma \sphinxparam{\DUrole{n}{prem}\DUrole{o}{=}\DUrole{default_value}{False}}}{}
\pysigstopsignatures
\sphinxAtStartPar
Escala la distribucion de slip a la magnitud de momento deseada de entrada.
crea una matriz de slip uniforme con magnitud de momento igual a la primitiva
para normalizar el slip. Despues escala con matriz uniforme con magnitud de momento
deseada
Entradas:
Mw: magnitud de momento deseada
slip: distribucion “primitiva” de slip
prof: matriz de profundidades de la interfaz
lons: longitudes de las subfallas
lats: latitudes de las subfallas
observacion: los 4 ultimos argumentos son para calcular la magnitud de momento 
con la funcion “magnitud\_momento” definida anteriormente

\end{fulllineitems}

\index{estima\_rigidez() (in module main)@\spxentry{estima\_rigidez()}\spxextra{in module main}}

\begin{fulllineitems}
\phantomsection\label{\detokenize{main:main.estima_rigidez}}
\pysigstartsignatures
\pysiglinewithargsret{\sphinxcode{\sphinxupquote{main.}}\sphinxbfcode{\sphinxupquote{estima\_rigidez}}}{\sphinxparam{\DUrole{n}{largo\_falla}}\sphinxparamcomma \sphinxparam{\DUrole{n}{tau}}}{}
\pysigstopsignatures
\sphinxAtStartPar
Estima la rigidez de la falla dado el tiempo de ruptura y el largo de la falla
segun la relacion dada en Bilek y Lay, 1999 mu = densidad*largo\_falla**2/(0.8**2*tau**2).
entradas: 
largo\_falla: largo de la falla en metros
tau: tiempo de duracion de la ruptura

\end{fulllineitems}

\index{estima\_rigidez\_prem() (in module main)@\spxentry{estima\_rigidez\_prem()}\spxextra{in module main}}

\begin{fulllineitems}
\phantomsection\label{\detokenize{main:main.estima_rigidez_prem}}
\pysigstartsignatures
\pysiglinewithargsret{\sphinxcode{\sphinxupquote{main.}}\sphinxbfcode{\sphinxupquote{estima\_rigidez\_prem}}}{\sphinxparam{\DUrole{n}{profs}}}{}
\pysigstopsignatures
\sphinxAtStartPar
Estima la rigidez de cada subfalla en funcion de su profunidad a partir de los 
datos de velocidad de onda S y densdidad del modelo PREM, dada Vs=sqrt(mu/rho)
Entradas:
profs: profundidades de las subfallas

\end{fulllineitems}

\index{hypocenter() (in module main)@\spxentry{hypocenter()}\spxextra{in module main}}

\begin{fulllineitems}
\phantomsection\label{\detokenize{main:main.hypocenter}}
\pysigstartsignatures
\pysiglinewithargsret{\sphinxcode{\sphinxupquote{main.}}\sphinxbfcode{\sphinxupquote{hypocenter}}}{\sphinxparam{\DUrole{n}{lons}}\sphinxparamcomma \sphinxparam{\DUrole{n}{lats}}\sphinxparamcomma \sphinxparam{\DUrole{n}{depth}}\sphinxparamcomma \sphinxparam{\DUrole{n}{length}}\sphinxparamcomma \sphinxparam{\DUrole{n}{width}}}{}
\pysigstopsignatures
\end{fulllineitems}

\index{interp\_slabtofault() (in module main)@\spxentry{interp\_slabtofault()}\spxextra{in module main}}

\begin{fulllineitems}
\phantomsection\label{\detokenize{main:main.interp_slabtofault}}
\pysigstartsignatures
\pysiglinewithargsret{\sphinxcode{\sphinxupquote{main.}}\sphinxbfcode{\sphinxupquote{interp\_slabtofault}}}{\sphinxparam{\DUrole{n}{lons}}\sphinxparamcomma \sphinxparam{\DUrole{n}{lats}}\sphinxparamcomma \sphinxparam{\DUrole{n}{nx}}\sphinxparamcomma \sphinxparam{\DUrole{n}{ny}}\sphinxparamcomma \sphinxparam{\DUrole{n}{slabdep}}\sphinxparamcomma \sphinxparam{\DUrole{n}{slabdip}}\sphinxparamcomma \sphinxparam{\DUrole{n}{slabstrike}}\sphinxparamcomma \sphinxparam{\DUrole{n}{slabrake}}}{}
\pysigstopsignatures
\end{fulllineitems}

\index{km2deg() (in module main)@\spxentry{km2deg()}\spxextra{in module main}}

\begin{fulllineitems}
\phantomsection\label{\detokenize{main:main.km2deg}}
\pysigstartsignatures
\pysiglinewithargsret{\sphinxcode{\sphinxupquote{main.}}\sphinxbfcode{\sphinxupquote{km2deg}}}{\sphinxparam{\DUrole{n}{km}}}{}
\pysigstopsignatures
\end{fulllineitems}

\index{load\_PREM() (in module main)@\spxentry{load\_PREM()}\spxextra{in module main}}

\begin{fulllineitems}
\phantomsection\label{\detokenize{main:main.load_PREM}}
\pysigstartsignatures
\pysiglinewithargsret{\sphinxcode{\sphinxupquote{main.}}\sphinxbfcode{\sphinxupquote{load\_PREM}}}{}{}
\pysigstopsignatures
\sphinxAtStartPar
carga los datos de velocidad de onda de corte del modelo prem
directorio: directorio donde esta el archivo del modelo slab2

\end{fulllineitems}

\index{load\_files\_slab2() (in module main)@\spxentry{load\_files\_slab2()}\spxextra{in module main}}

\begin{fulllineitems}
\phantomsection\label{\detokenize{main:main.load_files_slab2}}
\pysigstartsignatures
\pysiglinewithargsret{\sphinxcode{\sphinxupquote{main.}}\sphinxbfcode{\sphinxupquote{load\_files\_slab2}}}{\sphinxparam{\DUrole{n}{zone}\DUrole{o}{=}\DUrole{default_value}{\textquotesingle{}south\_america\textquotesingle{}}}\sphinxparamcomma \sphinxparam{\DUrole{n}{rake}\DUrole{o}{=}\DUrole{default_value}{True}}}{}
\pysigstopsignatures
\sphinxAtStartPar
Load data from the hayes 2018 slab2.0 model.

\sphinxAtStartPar
Inputs:
zone : str
\begin{quote}

\sphinxAtStartPar
Subduction zone to fech the model.
Available zones:
\begin{itemize}
\item {} 
\sphinxAtStartPar
\sphinxcode{\sphinxupquote{alaska}}: Alaska

\item {} 
\sphinxAtStartPar
\sphinxcode{\sphinxupquote{calabria}}: Calabria

\item {} 
\sphinxAtStartPar
\sphinxcode{\sphinxupquote{caribbean}}: Caribbean

\item {} 
\sphinxAtStartPar
\sphinxcode{\sphinxupquote{cascadia}}: Cascadia

\item {} 
\sphinxAtStartPar
\sphinxcode{\sphinxupquote{central\_america}}: Central America

\item {} 
\sphinxAtStartPar
\sphinxcode{\sphinxupquote{cotabalo}}: Cotabalo

\item {} 
\sphinxAtStartPar
\sphinxcode{\sphinxupquote{halmahera}}: Halmahera

\item {} 
\sphinxAtStartPar
\sphinxcode{\sphinxupquote{hellenic}}: Hellenic Arc

\item {} 
\sphinxAtStartPar
\sphinxcode{\sphinxupquote{himalaya}}: Himalaya

\item {} 
\sphinxAtStartPar
\sphinxcode{\sphinxupquote{hindu\_kush}}: Hindu Kush

\item {} 
\sphinxAtStartPar
\sphinxcode{\sphinxupquote{izu\_bonin}}: Izu\sphinxhyphen{}Bonin

\item {} 
\sphinxAtStartPar
\sphinxcode{\sphinxupquote{kamchatka}}: Kamchatka\sphinxhyphen{}Kuril Islands\sphinxhyphen{}Japan

\item {} 
\sphinxAtStartPar
\sphinxcode{\sphinxupquote{kermadec}}: Kermadec

\item {} 
\sphinxAtStartPar
\sphinxcode{\sphinxupquote{makran}}: Makran

\item {} 
\sphinxAtStartPar
\sphinxcode{\sphinxupquote{manila\_trench}}: Manila Trench

\item {} 
\sphinxAtStartPar
\sphinxcode{\sphinxupquote{muertos\_trough}}: Muertos Trough

\item {} 
\sphinxAtStartPar
\sphinxcode{\sphinxupquote{new\_guinea}}: New Guinea

\item {} 
\sphinxAtStartPar
\sphinxcode{\sphinxupquote{pamir}}: Pamir

\item {} 
\sphinxAtStartPar
\sphinxcode{\sphinxupquote{philippines}}: Philippines

\item {} 
\sphinxAtStartPar
\sphinxcode{\sphinxupquote{puysegur}}: Puysegur

\item {} 
\sphinxAtStartPar
\sphinxcode{\sphinxupquote{ryukyu}}: Ryukyu

\item {} 
\sphinxAtStartPar
\sphinxcode{\sphinxupquote{scotia\_sea}}: Scotia Sea

\item {} 
\sphinxAtStartPar
\sphinxcode{\sphinxupquote{solomon\_islands}}: Solomon Islands

\item {} 
\sphinxAtStartPar
\sphinxcode{\sphinxupquote{south\_america}}: South America

\item {} 
\sphinxAtStartPar
\sphinxcode{\sphinxupquote{sulawesi}}: Sulawesi

\item {} 
\sphinxAtStartPar
\sphinxcode{\sphinxupquote{sumatra\_java}}: Sumatra\sphinxhyphen{}Java

\item {} 
\sphinxAtStartPar
\sphinxcode{\sphinxupquote{vanuatu}}: Vanuatu

\end{itemize}
\end{quote}

\sphinxAtStartPar
If you have and wanna load rake file does nothing
if not, invoke function like load\_file\_slab2(zone,False)

\end{fulllineitems}

\index{load\_trench() (in module main)@\spxentry{load\_trench()}\spxextra{in module main}}

\begin{fulllineitems}
\phantomsection\label{\detokenize{main:main.load_trench}}
\pysigstartsignatures
\pysiglinewithargsret{\sphinxcode{\sphinxupquote{main.}}\sphinxbfcode{\sphinxupquote{load\_trench}}}{\sphinxparam{\DUrole{n}{trench\_file}}}{}
\pysigstopsignatures
\sphinxAtStartPar
upload file with the location of the south american plate trench
Input: trench file

\end{fulllineitems}

\index{magnitud\_momento() (in module main)@\spxentry{magnitud\_momento()}\spxextra{in module main}}

\begin{fulllineitems}
\phantomsection\label{\detokenize{main:main.magnitud_momento}}
\pysigstartsignatures
\pysiglinewithargsret{\sphinxcode{\sphinxupquote{main.}}\sphinxbfcode{\sphinxupquote{magnitud\_momento}}}{\sphinxparam{\DUrole{n}{slip}}\sphinxparamcomma \sphinxparam{\DUrole{n}{prof}}\sphinxparamcomma \sphinxparam{\DUrole{n}{lons}}\sphinxparamcomma \sphinxparam{\DUrole{n}{lats}}}{}
\pysigstopsignatures
\sphinxAtStartPar
Calcula la magnitud de momento Mw de la distribucion de slip con la ecuacion
Mw = 2/3log(Mo)\sphinxhyphen{}6.05
donde Mo = mu * area de ruptura * slip
Entradas: 
slip: corresponde a la distribucion de slip 
prof: matriz con profundidades
lons: matriz con longitudes
lats: matriz con latitudes

\end{fulllineitems}

\index{magnitud\_momento\_PREM() (in module main)@\spxentry{magnitud\_momento\_PREM()}\spxextra{in module main}}

\begin{fulllineitems}
\phantomsection\label{\detokenize{main:main.magnitud_momento_PREM}}
\pysigstartsignatures
\pysiglinewithargsret{\sphinxcode{\sphinxupquote{main.}}\sphinxbfcode{\sphinxupquote{magnitud\_momento\_PREM}}}{\sphinxparam{\DUrole{n}{slip}}\sphinxparamcomma \sphinxparam{\DUrole{n}{prof}}\sphinxparamcomma \sphinxparam{\DUrole{n}{lons}}\sphinxparamcomma \sphinxparam{\DUrole{n}{lats}}}{}
\pysigstopsignatures
\sphinxAtStartPar
Calcula la magnitud de momento Mw de la distribucion de slip con la ecuacion
Mw = 2/3log(Mo)\sphinxhyphen{}6.05. Calcula la rigidez individual de cada subfalla dependiendo de
su profundidad, a partir de los datos de velocidad de onda de corte y densidad del 
modelo PREM
donde Mo = mu * area de ruptura * slip
Entradas: 
slip: corresponde a la distribucion de slip 
prof: matriz con profundidades
lons: matriz con longitudes
lats: matriz con latitudes

\end{fulllineitems}

\index{make\_fault\_alongtrench() (in module main)@\spxentry{make\_fault\_alongtrench()}\spxextra{in module main}}

\begin{fulllineitems}
\phantomsection\label{\detokenize{main:main.make_fault_alongtrench}}
\pysigstartsignatures
\pysiglinewithargsret{\sphinxcode{\sphinxupquote{main.}}\sphinxbfcode{\sphinxupquote{make\_fault\_alongtrench}}}{\sphinxparam{\DUrole{n}{lons\_trench}}\sphinxparamcomma \sphinxparam{\DUrole{n}{lats\_trench}}\sphinxparamcomma \sphinxparam{\DUrole{n}{northlat}}\sphinxparamcomma \sphinxparam{\DUrole{n}{nx}}\sphinxparamcomma \sphinxparam{\DUrole{n}{ny}}\sphinxparamcomma \sphinxparam{\DUrole{n}{width}}\sphinxparamcomma \sphinxparam{\DUrole{n}{length}}}{}
\pysigstopsignatures
\sphinxAtStartPar
Function that creates a fault by calculating the latitudes and longitudes of various subfaults, considering the trench as a boundary.
Inputs:

\end{fulllineitems}

\index{make\_fault\_alongtrench\_optimized() (in module main)@\spxentry{make\_fault\_alongtrench\_optimized()}\spxextra{in module main}}

\begin{fulllineitems}
\phantomsection\label{\detokenize{main:main.make_fault_alongtrench_optimized}}
\pysigstartsignatures
\pysiglinewithargsret{\sphinxcode{\sphinxupquote{main.}}\sphinxbfcode{\sphinxupquote{make\_fault\_alongtrench\_optimized}}}{\sphinxparam{\DUrole{n}{lons\_trench}}\sphinxparamcomma \sphinxparam{\DUrole{n}{lats\_trench}}\sphinxparamcomma \sphinxparam{\DUrole{n}{northlat}}\sphinxparamcomma \sphinxparam{\DUrole{n}{nx}}\sphinxparamcomma \sphinxparam{\DUrole{n}{ny}}\sphinxparamcomma \sphinxparam{\DUrole{n}{width}}\sphinxparamcomma \sphinxparam{\DUrole{n}{length}}}{}
\pysigstopsignatures
\end{fulllineitems}

\index{matriz\_covarianza() (in module main)@\spxentry{matriz\_covarianza()}\spxextra{in module main}}

\begin{fulllineitems}
\phantomsection\label{\detokenize{main:main.matriz_covarianza}}
\pysigstartsignatures
\pysiglinewithargsret{\sphinxcode{\sphinxupquote{main.}}\sphinxbfcode{\sphinxupquote{matriz\_covarianza}}}{\sphinxparam{\DUrole{n}{dip}}\sphinxparamcomma \sphinxparam{\DUrole{n}{prof}}\sphinxparamcomma \sphinxparam{\DUrole{n}{lons}}\sphinxparamcomma \sphinxparam{\DUrole{n}{lats}}}{}
\pysigstopsignatures
\sphinxAtStartPar
calcula la matriz de correlacion:
Cij = exp(\sphinxhyphen{}(dstrike(i,j)/rstrike)\sphinxhyphen{}(ddip(i,j)/rdip)
donde:
dstrike y ddip son estimados de las distancias entre las subfallas i \& j
ddip = dprof/sin(dip)
d es la distancia euclidiana entre dos puntos aproximada segun sus lats y lons
rstrike y rdip son los largos de correlacion en cada direccion 
Entradas: 
dip: matriz con los dips de cada punto
prof: matriz con las profundidades de cada punto
lons: matriz con longitud de cada subfalla, necesaria para el calculo de las distancias entre estas
lats: matriz con latitud de cada subfalla, necesaria para el calculo de las distancias entre estas
Observacion: en esta primera instancia se considera la aproximacion  1 grado approx 111.12 km (generalizar para fallas mas grandes)
Observacion 2: aproximacion para calculo de distancia corregida con calculo geodesico de distancia (aumenta el costo computacional bastante)

\end{fulllineitems}

\index{matriz\_covarianza\_optimized() (in module main)@\spxentry{matriz\_covarianza\_optimized()}\spxextra{in module main}}

\begin{fulllineitems}
\phantomsection\label{\detokenize{main:main.matriz_covarianza_optimized}}
\pysigstartsignatures
\pysiglinewithargsret{\sphinxcode{\sphinxupquote{main.}}\sphinxbfcode{\sphinxupquote{matriz\_covarianza\_optimized}}}{\sphinxparam{\DUrole{n}{dip}}\sphinxparamcomma \sphinxparam{\DUrole{n}{prof}}\sphinxparamcomma \sphinxparam{\DUrole{n}{lons}}\sphinxparamcomma \sphinxparam{\DUrole{n}{lats}}\sphinxparamcomma \sphinxparam{\DUrole{n}{largo\_falla}}\sphinxparamcomma \sphinxparam{\DUrole{n}{ancho\_falla}}\sphinxparamcomma \sphinxparam{\DUrole{n}{alpha}\DUrole{o}{=}\DUrole{default_value}{0.5}}}{}
\pysigstopsignatures
\sphinxAtStartPar
calcula la matriz de correlacion:
Cij = exp(\sphinxhyphen{}(dstrike(i,j)/rstrike)\sphinxhyphen{}(ddip(i,j)/rdip)
donde:
dstrike y ddip son estimados de las distancias entre las subfallas i \& j
ddip = dprof/sin(dip)
d es la distancia euclidiana entre dos puntos aproximada segun sus lats y lons
rstrike y rdip son los largos de correlacion en cada direccion 
Entradas: 
dip: matriz con los dips de cada punto
prof: matriz con las profundidades de cada punto
lons: matriz con longitud de cada subfalla, necesaria para el calculo de las distancias entre estas
lats: matriz con latitud de cada subfalla, necesaria para el calculo de las distancias entre estas
Observacion: en esta primera instancia se considera la aproximacion  1 grado approx 111.12 km (generalizar para fallas mas grandes)
Observacion 2: aproximacion para calculo de distancia corregida con calculo geodesico de distancia (aumenta el costo computacional bastante)
Obeservacion 3: el ancho y largo de la falla debe ser en metros

\end{fulllineitems}

\index{matriz\_medias() (in module main)@\spxentry{matriz\_medias()}\spxextra{in module main}}

\begin{fulllineitems}
\phantomsection\label{\detokenize{main:main.matriz_medias}}
\pysigstartsignatures
\pysiglinewithargsret{\sphinxcode{\sphinxupquote{main.}}\sphinxbfcode{\sphinxupquote{matriz\_medias}}}{\sphinxparam{\DUrole{n}{media}}\sphinxparamcomma \sphinxparam{\DUrole{n}{prof}}}{}
\pysigstopsignatures
\sphinxAtStartPar
calcula las medias mu\_i para cada subfalla segun mu\_i = log(mu*tau\_i)\sphinxhyphen{}1/2log(alpha**2+1)
Entradas:
media: valor promedio alrededor del que se desea que se centre el slip
prof: matriz con profundidades de cada subfalla, necesaria para el taper

\end{fulllineitems}

\index{media\_slip() (in module main)@\spxentry{media\_slip()}\spxextra{in module main}}

\begin{fulllineitems}
\phantomsection\label{\detokenize{main:main.media_slip}}
\pysigstartsignatures
\pysiglinewithargsret{\sphinxcode{\sphinxupquote{main.}}\sphinxbfcode{\sphinxupquote{media\_slip}}}{\sphinxparam{\DUrole{n}{Mw}}\sphinxparamcomma \sphinxparam{\DUrole{n}{largo}}\sphinxparamcomma \sphinxparam{\DUrole{n}{ancho}}\sphinxparamcomma \sphinxparam{\DUrole{n}{prof}}}{}
\pysigstopsignatures
\sphinxAtStartPar
Largo, ancho y profundidad deben estar en metros

\end{fulllineitems}

\index{nearest\_value() (in module main)@\spxentry{nearest\_value()}\spxextra{in module main}}

\begin{fulllineitems}
\phantomsection\label{\detokenize{main:main.nearest_value}}
\pysigstartsignatures
\pysiglinewithargsret{\sphinxcode{\sphinxupquote{main.}}\sphinxbfcode{\sphinxupquote{nearest\_value}}}{\sphinxparam{\DUrole{n}{value}}\sphinxparamcomma \sphinxparam{\DUrole{n}{matrix}}}{}
\pysigstopsignatures
\end{fulllineitems}

\index{plot\_3d() (in module main)@\spxentry{plot\_3d()}\spxextra{in module main}}

\begin{fulllineitems}
\phantomsection\label{\detokenize{main:main.plot_3d}}
\pysigstartsignatures
\pysiglinewithargsret{\sphinxcode{\sphinxupquote{main.}}\sphinxbfcode{\sphinxupquote{plot\_3d}}}{\sphinxparam{\DUrole{n}{X\_array}}\sphinxparamcomma \sphinxparam{\DUrole{n}{Y\_array}}\sphinxparamcomma \sphinxparam{\DUrole{n}{depth}}\sphinxparamcomma \sphinxparam{\DUrole{n}{Slip}}\sphinxparamcomma \sphinxparam{\DUrole{n}{filename}\DUrole{o}{=}\DUrole{default_value}{None}}}{}
\pysigstopsignatures
\end{fulllineitems}

\index{plot\_grid() (in module main)@\spxentry{plot\_grid()}\spxextra{in module main}}

\begin{fulllineitems}
\phantomsection\label{\detokenize{main:main.plot_grid}}
\pysigstartsignatures
\pysiglinewithargsret{\sphinxcode{\sphinxupquote{main.}}\sphinxbfcode{\sphinxupquote{plot\_grid}}}{}{}
\pysigstopsignatures
\end{fulllineitems}

\index{plot\_slip() (in module main)@\spxentry{plot\_slip()}\spxextra{in module main}}

\begin{fulllineitems}
\phantomsection\label{\detokenize{main:main.plot_slip}}
\pysigstartsignatures
\pysiglinewithargsret{\sphinxcode{\sphinxupquote{main.}}\sphinxbfcode{\sphinxupquote{plot\_slip}}}{\sphinxparam{\DUrole{n}{X\_grid}}\sphinxparamcomma \sphinxparam{\DUrole{n}{Y\_grid}}\sphinxparamcomma \sphinxparam{\DUrole{n}{lonfosa}}\sphinxparamcomma \sphinxparam{\DUrole{n}{latfosa}}\sphinxparamcomma \sphinxparam{\DUrole{n}{Slip}}\sphinxparamcomma \sphinxparam{\DUrole{n}{filename}}\sphinxparamcomma \sphinxparam{\DUrole{n}{show}\DUrole{o}{=}\DUrole{default_value}{False}}\sphinxparamcomma \sphinxparam{\DUrole{n}{cmap}\DUrole{o}{=}\DUrole{default_value}{\textquotesingle{}rainbow\textquotesingle{}}}}{}
\pysigstopsignatures
\end{fulllineitems}

\index{plot\_slip\_css() (in module main)@\spxentry{plot\_slip\_css()}\spxextra{in module main}}

\begin{fulllineitems}
\phantomsection\label{\detokenize{main:main.plot_slip_css}}
\pysigstartsignatures
\pysiglinewithargsret{\sphinxcode{\sphinxupquote{main.}}\sphinxbfcode{\sphinxupquote{plot\_slip\_css}}}{\sphinxparam{\DUrole{n}{region}}\sphinxparamcomma \sphinxparam{\DUrole{n}{lons}}\sphinxparamcomma \sphinxparam{\DUrole{n}{lats}}\sphinxparamcomma \sphinxparam{\DUrole{n}{lonfosa}}\sphinxparamcomma \sphinxparam{\DUrole{n}{latfosa}}\sphinxparamcomma \sphinxparam{\DUrole{n}{Slip}}\sphinxparamcomma \sphinxparam{\DUrole{n}{cmap}\DUrole{o}{=}\DUrole{default_value}{\textquotesingle{}rainbow\textquotesingle{}}}}{}
\pysigstopsignatures
\end{fulllineitems}

\index{plot\_slip\_gmt() (in module main)@\spxentry{plot\_slip\_gmt()}\spxextra{in module main}}

\begin{fulllineitems}
\phantomsection\label{\detokenize{main:main.plot_slip_gmt}}
\pysigstartsignatures
\pysiglinewithargsret{\sphinxcode{\sphinxupquote{main.}}\sphinxbfcode{\sphinxupquote{plot\_slip\_gmt}}}{\sphinxparam{\DUrole{n}{region}}\sphinxparamcomma \sphinxparam{\DUrole{n}{X\_grid}}\sphinxparamcomma \sphinxparam{\DUrole{n}{Y\_grid}}\sphinxparamcomma \sphinxparam{\DUrole{n}{lonfosa}}\sphinxparamcomma \sphinxparam{\DUrole{n}{latfosa}}\sphinxparamcomma \sphinxparam{\DUrole{n}{Slip}}\sphinxparamcomma \sphinxparam{\DUrole{n}{dx}}\sphinxparamcomma \sphinxparam{\DUrole{n}{dy}}\sphinxparamcomma \sphinxparam{\DUrole{n}{archivo\_salida}}}{}
\pysigstopsignatures
\end{fulllineitems}

\index{sigmoid() (in module main)@\spxentry{sigmoid()}\spxextra{in module main}}

\begin{fulllineitems}
\phantomsection\label{\detokenize{main:main.sigmoid}}
\pysigstartsignatures
\pysiglinewithargsret{\sphinxcode{\sphinxupquote{main.}}\sphinxbfcode{\sphinxupquote{sigmoid}}}{\sphinxparam{\DUrole{n}{x}}}{}
\pysigstopsignatures
\sphinxAtStartPar
Funcion auxiliar para crear ventana sigmoidal para el taper

\end{fulllineitems}

\index{taper\_LeVeque() (in module main)@\spxentry{taper\_LeVeque()}\spxextra{in module main}}

\begin{fulllineitems}
\phantomsection\label{\detokenize{main:main.taper_LeVeque}}
\pysigstartsignatures
\pysiglinewithargsret{\sphinxcode{\sphinxupquote{main.}}\sphinxbfcode{\sphinxupquote{taper\_LeVeque}}}{\sphinxparam{\DUrole{n}{d}}}{}
\pysigstopsignatures
\sphinxAtStartPar
realiza el taper propuesto por leveque et al., 2016
t(d)=1\sphinxhyphen{}exp(\sphinxhyphen{}20|d\sphinxhyphen{}dmax|/dmax) 
donde dmax es la profunidad maxima y de la falla y d es la profundidad
de cada subfalla
Entrada:
d: matriz con profundidades

\end{fulllineitems}

\index{taper\_borders() (in module main)@\spxentry{taper\_borders()}\spxextra{in module main}}

\begin{fulllineitems}
\phantomsection\label{\detokenize{main:main.taper_borders}}
\pysigstartsignatures
\pysiglinewithargsret{\sphinxcode{\sphinxupquote{main.}}\sphinxbfcode{\sphinxupquote{taper\_borders}}}{\sphinxparam{\DUrole{n}{Slip}}\sphinxparamcomma \sphinxparam{\DUrole{n}{taperfunc=\textless{}function hanning\textgreater{}}}}{}
\pysigstopsignatures
\end{fulllineitems}

\index{taper\_except\_trench() (in module main)@\spxentry{taper\_except\_trench()}\spxextra{in module main}}

\begin{fulllineitems}
\phantomsection\label{\detokenize{main:main.taper_except_trench}}
\pysigstartsignatures
\pysiglinewithargsret{\sphinxcode{\sphinxupquote{main.}}\sphinxbfcode{\sphinxupquote{taper\_except\_trench}}}{\sphinxparam{\DUrole{n}{Slip}}\sphinxparamcomma \sphinxparam{\DUrole{n}{taperfunc=\textless{}function hanning\textgreater{}}}}{}
\pysigstopsignatures
\end{fulllineitems}

\index{taper\_except\_trench\_tukey() (in module main)@\spxentry{taper\_except\_trench\_tukey()}\spxextra{in module main}}

\begin{fulllineitems}
\phantomsection\label{\detokenize{main:main.taper_except_trench_tukey}}
\pysigstartsignatures
\pysiglinewithargsret{\sphinxcode{\sphinxupquote{main.}}\sphinxbfcode{\sphinxupquote{taper\_except\_trench\_tukey}}}{\sphinxparam{\DUrole{n}{Slip}}\sphinxparamcomma \sphinxparam{\DUrole{n}{dip\_taperfunc=\textless{}function tukey\_window\textgreater{}}}\sphinxparamcomma \sphinxparam{\DUrole{n}{strike\_taperfunc=\textless{}function tukey\textgreater{}}}\sphinxparamcomma \sphinxparam{\DUrole{n}{alpha\_dip=0.5}}\sphinxparamcomma \sphinxparam{\DUrole{n}{alpha\_strike=0.5}}}{}
\pysigstopsignatures
\end{fulllineitems}

\index{taper\_slip\_fosa() (in module main)@\spxentry{taper\_slip\_fosa()}\spxextra{in module main}}

\begin{fulllineitems}
\phantomsection\label{\detokenize{main:main.taper_slip_fosa}}
\pysigstartsignatures
\pysiglinewithargsret{\sphinxcode{\sphinxupquote{main.}}\sphinxbfcode{\sphinxupquote{taper\_slip\_fosa}}}{\sphinxparam{\DUrole{n}{Slip}}\sphinxparamcomma \sphinxparam{\DUrole{n}{ventana}}}{}
\pysigstopsignatures
\sphinxAtStartPar
Aplica un taper para atenuar el slip en las cercanias de la fosa
Entradas:
Slip: array con los valores de slip por subfalla
ventana: ventana a usarse para aplicar el taper

\end{fulllineitems}

\index{tau\_ruptura() (in module main)@\spxentry{tau\_ruptura()}\spxextra{in module main}}

\begin{fulllineitems}
\phantomsection\label{\detokenize{main:main.tau_ruptura}}
\pysigstartsignatures
\pysiglinewithargsret{\sphinxcode{\sphinxupquote{main.}}\sphinxbfcode{\sphinxupquote{tau\_ruptura}}}{\sphinxparam{\DUrole{n}{largo\_falla}}\sphinxparamcomma \sphinxparam{\DUrole{n}{beta}\DUrole{o}{=}\DUrole{default_value}{2500}}}{}
\pysigstopsignatures
\sphinxAtStartPar
Calcula el tiempo de ruptura de la falla dada a partir del largo de falla, considerando
la velocidad de ruptura como 0.8 de la velocidad de propagacion de onda S y 
una velocidad de cizalle de 2.5 km/s utilizando la aproximacion dada en Bilek y Lay, 1999
tau approx L/(0.8beta) 
Entradas: 
largo\_falla: largo de la falla en metros dado por la funcion que crea la falla
beta: velocidad de la onda de cizalle en metros/segundo

\end{fulllineitems}

\index{tau\_ruptura\_prem() (in module main)@\spxentry{tau\_ruptura\_prem()}\spxextra{in module main}}

\begin{fulllineitems}
\phantomsection\label{\detokenize{main:main.tau_ruptura_prem}}
\pysigstartsignatures
\pysiglinewithargsret{\sphinxcode{\sphinxupquote{main.}}\sphinxbfcode{\sphinxupquote{tau\_ruptura\_prem}}}{\sphinxparam{\DUrole{n}{largo\_falla}}\sphinxparamcomma \sphinxparam{\DUrole{n}{prof\_media}}\sphinxparamcomma \sphinxparam{\DUrole{n}{prof\_prem}}\sphinxparamcomma \sphinxparam{\DUrole{n}{vs\_prem}}}{}
\pysigstopsignatures
\sphinxAtStartPar
Calcula el tiempo de ruptura de la falla dada a partir del largo de falla, considerando
la velocidad de ruptura como 0.8 de la velocidad de propagacion de onda S y 
una velocidad de cizalle dependiente de la profundidad media de la falla
y la velocidad vs asociada a esta en el modelo prem, utilizando la aproximacion dada en Bilek y Lay, 1999
tau approx L/(0.8beta) 
Entradas: 
largo\_falla: largo de la falla en metros dado por la funcion que crea la falla
prof\_media: profundidad media de la falla
prof\_prem: profundidades del modelo prem
vs: modelo de velocidades del modelo prem asociadas a las profundidades de prof\_prem

\end{fulllineitems}

\index{tukey\_window() (in module main)@\spxentry{tukey\_window()}\spxextra{in module main}}

\begin{fulllineitems}
\phantomsection\label{\detokenize{main:main.tukey_window}}
\pysigstartsignatures
\pysiglinewithargsret{\sphinxcode{\sphinxupquote{main.}}\sphinxbfcode{\sphinxupquote{tukey\_window}}}{\sphinxparam{\DUrole{n}{N}}\sphinxparamcomma \sphinxparam{\DUrole{n}{alpha}\DUrole{o}{=}\DUrole{default_value}{0.5}}}{}
\pysigstopsignatures
\sphinxAtStartPar
Genera una ventana Tukey.Se ha
ajustado los valores de la parte de Hann para que no alcancen completamente cero al comienzo y al final

\sphinxAtStartPar
Parámetros:
\sphinxhyphen{} N: Longitud de la ventana.
\sphinxhyphen{} alpha: Parámetro de apertura (0 para una ventana rectangular, 1 para una ventana Hann).

\sphinxAtStartPar
Retorna:
\sphinxhyphen{} Ventana Tukey.

\end{fulllineitems}

\index{ventana\_taper\_slip\_fosa() (in module main)@\spxentry{ventana\_taper\_slip\_fosa()}\spxextra{in module main}}

\begin{fulllineitems}
\phantomsection\label{\detokenize{main:main.ventana_taper_slip_fosa}}
\pysigstartsignatures
\pysiglinewithargsret{\sphinxcode{\sphinxupquote{main.}}\sphinxbfcode{\sphinxupquote{ventana\_taper\_slip\_fosa}}}{\sphinxparam{\DUrole{n}{Slip}}\sphinxparamcomma \sphinxparam{\DUrole{n}{lon}}\sphinxparamcomma \sphinxparam{\DUrole{n}{lat}}\sphinxparamcomma \sphinxparam{\DUrole{n}{ventana\_flag}}}{}
\pysigstopsignatures
\sphinxAtStartPar
Crea una ventana para aplicar un taper al slip en cercanias de la fosa
Entradas:
Slip: array con los valores de slip por subfalla
lon: array de longitudes de las subfallas
lat: array con las latitudes de las subfallas

\end{fulllineitems}



\subsection{Submodules}
\label{\detokenize{main:submodules}}

\subsection{main.modcomcot module}
\label{\detokenize{main:module-main.modcomcot}}\label{\detokenize{main:main-modcomcot-module}}\index{module@\spxentry{module}!main.modcomcot@\spxentry{main.modcomcot}}\index{main.modcomcot@\spxentry{main.modcomcot}!module@\spxentry{module}}\index{make\_multifault() (in module main.modcomcot)@\spxentry{make\_multifault()}\spxextra{in module main.modcomcot}}

\begin{fulllineitems}
\phantomsection\label{\detokenize{main:main.modcomcot.make_multifault}}
\pysigstartsignatures
\pysiglinewithargsret{\sphinxcode{\sphinxupquote{main.modcomcot.}}\sphinxbfcode{\sphinxupquote{make\_multifault}}}{\sphinxparam{\DUrole{n}{file}}\sphinxparamcomma \sphinxparam{\DUrole{n}{fname}}\sphinxparamcomma \sphinxparam{\DUrole{n}{lat\_griddomain}}\sphinxparamcomma \sphinxparam{\DUrole{n}{lon\_griddomain}}}{}
\pysigstopsignatures
\end{fulllineitems}



\subsection{main.modfilters module}
\label{\detokenize{main:module-main.modfilters}}\label{\detokenize{main:main-modfilters-module}}\index{module@\spxentry{module}!main.modfilters@\spxentry{main.modfilters}}\index{main.modfilters@\spxentry{main.modfilters}!module@\spxentry{module}}\index{corr2\_coeff() (in module main.modfilters)@\spxentry{corr2\_coeff()}\spxextra{in module main.modfilters}}

\begin{fulllineitems}
\phantomsection\label{\detokenize{main:main.modfilters.corr2_coeff}}
\pysigstartsignatures
\pysiglinewithargsret{\sphinxcode{\sphinxupquote{main.modfilters.}}\sphinxbfcode{\sphinxupquote{corr2\_coeff}}}{\sphinxparam{\DUrole{n}{A}}\sphinxparamcomma \sphinxparam{\DUrole{n}{B}}}{}
\pysigstopsignatures
\end{fulllineitems}

\index{couplingfilter() (in module main.modfilters)@\spxentry{couplingfilter()}\spxextra{in module main.modfilters}}

\begin{fulllineitems}
\phantomsection\label{\detokenize{main:main.modfilters.couplingfilter}}
\pysigstartsignatures
\pysiglinewithargsret{\sphinxcode{\sphinxupquote{main.modfilters.}}\sphinxbfcode{\sphinxupquote{couplingfilter}}}{\sphinxparam{\DUrole{n}{X\_grid}}\sphinxparamcomma \sphinxparam{\DUrole{n}{Y\_grid}}\sphinxparamcomma \sphinxparam{\DUrole{n}{Slip}}\sphinxparamcomma \sphinxparam{\DUrole{n}{couplingfilename}}\sphinxparamcomma \sphinxparam{\DUrole{n}{lonfosa}}\sphinxparamcomma \sphinxparam{\DUrole{n}{latfosa}}}{}
\pysigstopsignatures
\end{fulllineitems}

\index{depthfilter() (in module main.modfilters)@\spxentry{depthfilter()}\spxextra{in module main.modfilters}}

\begin{fulllineitems}
\phantomsection\label{\detokenize{main:main.modfilters.depthfilter}}
\pysigstartsignatures
\pysiglinewithargsret{\sphinxcode{\sphinxupquote{main.modfilters.}}\sphinxbfcode{\sphinxupquote{depthfilter}}}{\sphinxparam{\DUrole{n}{Slip}}\sphinxparamcomma \sphinxparam{\DUrole{n}{depth}}\sphinxparamcomma \sphinxparam{\DUrole{n}{prctile}\DUrole{o}{=}\DUrole{default_value}{10}}}{}
\pysigstopsignatures
\end{fulllineitems}

\index{obtain\_border() (in module main.modfilters)@\spxentry{obtain\_border()}\spxextra{in module main.modfilters}}

\begin{fulllineitems}
\phantomsection\label{\detokenize{main:main.modfilters.obtain_border}}
\pysigstartsignatures
\pysiglinewithargsret{\sphinxcode{\sphinxupquote{main.modfilters.}}\sphinxbfcode{\sphinxupquote{obtain\_border}}}{\sphinxparam{\DUrole{n}{matrix}}}{}
\pysigstopsignatures
\end{fulllineitems}

\index{physical\_filter() (in module main.modfilters)@\spxentry{physical\_filter()}\spxextra{in module main.modfilters}}

\begin{fulllineitems}
\phantomsection\label{\detokenize{main:main.modfilters.physical_filter}}
\pysigstartsignatures
\pysiglinewithargsret{\sphinxcode{\sphinxupquote{main.modfilters.}}\sphinxbfcode{\sphinxupquote{physical\_filter}}}{\sphinxparam{\DUrole{n}{Slip}}}{}
\pysigstopsignatures
\end{fulllineitems}



\subsection{main.modokada module}
\label{\detokenize{main:main-modokada-module}}
\sphinxstepscope


\chapter{Overview}
\label{\detokenize{overview:overview}}\label{\detokenize{overview::doc}}
\sphinxstepscope


\chapter{Installing}
\label{\detokenize{installing:installing}}\label{\detokenize{installing::doc}}
\sphinxstepscope


\chapter{Changelog}
\label{\detokenize{changelog:changelog}}\label{\detokenize{changelog::doc}}

\section{Subsection}
\label{\detokenize{changelog:subsection}}

\chapter{Indices and tables}
\label{\detokenize{index:indices-and-tables}}\begin{itemize}
\item {} 
\sphinxAtStartPar
\DUrole{xref,std,std-ref}{genindex}

\item {} 
\sphinxAtStartPar
\DUrole{xref,std,std-ref}{modindex}

\item {} 
\sphinxAtStartPar
\DUrole{xref,std,std-ref}{search}

\end{itemize}


\renewcommand{\indexname}{Python Module Index}
\begin{sphinxtheindex}
\let\bigletter\sphinxstyleindexlettergroup
\bigletter{m}
\item\relax\sphinxstyleindexentry{main}\sphinxstyleindexpageref{main:\detokenize{module-main}}
\item\relax\sphinxstyleindexentry{main.modcomcot}\sphinxstyleindexpageref{main:\detokenize{module-main.modcomcot}}
\item\relax\sphinxstyleindexentry{main.modfilters}\sphinxstyleindexpageref{main:\detokenize{module-main.modfilters}}
\end{sphinxtheindex}

\renewcommand{\indexname}{Index}
\printindex
\end{document}